% ============================================================================
\section{GPAS}
\label{sec:theory}
% ============================================================================

\subsection{Noise in preconditioned coordinates}

The variance criterion is exact for a step with a fixed diagonal
preconditioner and provides a local allocation metric for optimizers with
diagonal preconditioning. Let $D$ be the pre-step diagonal preconditioner read
from the optimizer state:
\[
    D=
    \begin{cases}
        I, & \text{for SGD},\\
        \operatorname{diag}\!\bigl((\sqrt{\widehat v_{\rm pre}}+\epsilon)^{-1}\bigr),
        & \text{for an Adam-type optimizer},
    \end{cases}
\]
where $\widehat v_{\rm pre}$ is the bias-corrected second-moment state at the
start of the step. Noise in a coordinate affects the local update according to
how strongly $D$ amplifies it, so we measure noise after preconditioning and
hold $D$ fixed within a step.

The quantity to allocate is the variance of a single micro-batch gradient
around its task mean. Unless noted otherwise, $g_i$ and $e_i$ refer to the
top-$k$ estimator used for training. Under the working model that fresh
micro-batches are conditionally independent across draws and identically
distributed within each task, with a common task mean and covariance, define
task noise $e_i$ and the accumulated-gradient variance $V(m)$ by
\begin{equation}
    e_i=\mathbb E\bigl\|D(g_i-\mu_i)\bigr\|_2^2,
    \qquad
    V(m)=\mathbb E\bigl\|D(A-\mu)\bigr\|_2^2
        =\sum_{i=1}^M\frac{w_i^2e_i}{m_i}.
    \label{eq:batch_variance}
\end{equation}
Thus the allocation enters the preconditioned variance only through the
inverse-count terms. When $\mu$ is the gradient of the step-local surrogate,
then for a fixed student state, preconditioner, and step size, the allocation
also enters its one-step descent bound only through $V(m)$ (proof in
Appendix~\ref{app:descent}).

Minimizing this variance gives the count rule used by GPAS. Before enforcing
integer bounds, the solution is
\begin{equation}
    m_i\ \propto\ w_i\sqrt{e_i}.
    \label{eq:batch_gpas_allocation}
\end{equation}
This is Neyman allocation, the classical rule that samples each stratum in
proportion to its standard deviation. We choose the common proportionality
scale so that the clipped counts sum to $G$, iterating after any count reaches
a bound, and then round by largest remainder.

Preconditioning matters because unpreconditioned and preconditioned noise can
rank tasks differently. In a three-task example whose preconditioner reverses
the noise ranking, counts from unpreconditioned noise raise the preconditioned
variance to $2.3\times$ Uniform, whereas the preconditioned rule lowers it to
$0.68\times$ (details in Appendix~\ref{app:controlled}).

\subsection{Where the teacher--student gap enters}
\label{sec:gap_noise}

The log ratio at a prefix has two distinct roles. Its student-probability
weighted mean is the local reverse-KL gap, while its centered variation,
together with the norm and direction of
$\nabla_\theta\log\pi_\theta(y_t\mid x,y_{<t})$, determines how noisy a sampled
score-function term can be; the constant mean itself cancels against the
mean-zero score function. Teachers at different distances from the student can
therefore induce different within-position spreads, an empirical relation we
test rather than assume.

For a precise diagnostic, view scoring in two stages: first fix the prompts and
on-policy prefixes in a micro-batch, and then independently draw the token used
to estimate each fixed-prefix score. Let $g_{i,\mathrm{diag}}^a$ denote the
resulting diagnostic gradient for estimator
$a\in\{\mathrm{samp},\mathrm{top}k\}$, let $\bar g_i$ be their common
conditional mean given the fixed trace, and let
$\mu_i^{\mathrm{diag}}=\mathbb E[g_{i,\mathrm{diag}}^a]$. The law of total
variance gives
\[
    e_{i,\mathrm{diag}}^a
      =e_{i,\mathrm{diag}}^{\mathrm{resp}}
       +e_{i,\mathrm{diag}}^{\mathrm{tok},a},\qquad
    e_{i,\mathrm{diag}}^{\mathrm{resp}}
      =\mathbb E\bigl\|D(\bar g_i-\mu_i^{\mathrm{diag}})\bigr\|_2^2,
    \quad
    e_{i,\mathrm{diag}}^{\mathrm{tok},a}
      =\mathbb E\bigl\|D(g_{i,\mathrm{diag}}^a-\bar g_i)\bigr\|_2^2.
\]
Here response-level noise comes from which prompts and rollout prefixes are
observed, whereas token-level noise is the sum of the conditional
within-position contributions from which scoring token is drawn at each fixed
prefix (derivation in Appendix~\ref{app:allocation_proofs}). GPAS handles
response-level noise by averaging more examples, whereas the retained-token
component can be integrated directly at a fixed prefix.

The estimator in Section~\ref{sec:setting} replaces the random retained-set
contribution by its exact fixed-prefix expectation. Its tail correction remains
random, so the paired measurement in Section~\ref{sec:experiments} checks the
diagnostic variance change rather than presuming it. Online GPAS uses $e_i$
from the actual top-$k$ training gradients in Equation~\ref{eq:batch_variance};
the independent scoring-token protocol isolates the fixed-prefix mechanism,
and its diagnostic components need not equal that online $e_i$. We test whether
more distant teachers show a larger diagnostic token component and whether
integrating the retained contribution predicts a narrower online noise spread
and a changed allocation.

Open-MOPD uses the level of the gap as a budget signal and increases the weight
of large-gap domains \citep{gao2026openmopddiagnosingfixingcapability}. Under a
fixed objective, the log ratio already appears in the task gradient; its
additional token-to-token variability is estimation noise, which we address
with the estimator rather than by changing objective weights.

\subsection{Cost-aware mode}

Long-response tasks can take more time per micro-batch, so the allocation with
the least variance per step need not have the least variance per hour. Let
$\tau_i$ be the measured marginal time for the count-dependent rollout,
teacher scoring, and backward work of one task-$i$ micro-batch, and let $C$ be
the remaining fixed time for synchronization, the optimizer update, and
logging. At a fixed student state, under the measured additive time model, the
local time-normalized variance proxy is
\begin{equation}
    J_C(m)=\Bigl(C+\sum_{i=1}^M m_i\tau_i\Bigr)\,V(m).
    \label{eq:batch_cost_objective}
\end{equation}
GPAS in its cost-aware mode enumerates the feasible integer count vectors and
chooses the one with smallest $J_C$. In the unconstrained continuous relaxation
with no fixed time, the optimal count scales as
$w_i\sqrt{e_i/\tau_i}$; as fixed time dominates, it approaches the per-step
rule (derivation in Appendix~\ref{app:allocation_proofs}).

\subsection{Online estimates and implementation}

Every statistic for the next allocation comes from the step just completed.
Let $\bar g_i$ now denote the sample mean of task $i$'s observed micro-batch
gradients in that step. Its preconditioned sample variance is
\begin{equation}
    \widehat e_i=\frac{1}{m_i-1}\sum_{s=1}^{m_i}
    \bigl\|D(g_{i,s}-\bar g_i)\bigr\|_2^2.
    \label{eq:noise_estimate}
\end{equation}
Under this working model, it is an unbiased estimate of one-micro-batch
preconditioned noise, and the count lower bound guarantees that it exists every
step. Timing estimates are smoothed online because response lengths make
individual observations noisy.

To expose how much a count change can help, define the plug-in prediction
$\widehat V(m)=\sum_iw_i^2\widehat e_i/m_i$ and compare it with Uniform. The
logged variance ratio is
\begin{equation}
    H=\frac{\widehat V(m^{\rm unif})}{\widehat V(m)},
    \label{eq:allocation_variance_ratio}
\end{equation}
and is bounded above by the count limits. The first step uses Uniform because
no earlier statistics are available.

GPAS leaves the optimizer untouched: for SGD, $D$ is the identity and no
optimizer state is read; for Adam-type optimizers, $D$ is read from the
pre-step second-moment buffer. Changing the allocation changes the noise term
in an Adam-type second-moment observation, as does any change in batch
composition; SGD has no such state and is unaffected (details in
Appendix~\ref{app:moment_target_proof}). Computing $\widehat e_i$ requires one
parameter-sized mean buffer and elementwise reductions, but no additional
forward or backward pass. The top-$k$ estimator likewise reuses logits that
the student and teacher already compute. Algorithm~\ref{alg:batch_gpas}
summarizes one step.

\begin{algorithm}[H]
\caption{One optimizer step with GPAS.}
\label{alg:batch_gpas}
\small
\begin{minipage}{0.97\linewidth}
\begin{enumerate}
    \item Choose bounded integer counts $m$ from the previous step's
    $\widehat e_i$ and, in the cost-aware mode, $\tau_i$ and $C$; use Uniform
    when no statistics are available.
    \item For every task, draw fresh prompts and responses, obtain teacher
    scores for the student top-$k$ tokens and for the sampled token when it is
    outside that set, and accumulate each micro-batch gradient with weight
    $w_i/m_i$.
    \item Read the pre-update $D$ from the optimizer state (the identity for
    SGD); from the same gradients and trace, compute $\widehat e_i$, update
    timing estimates, and log the diagnostic quantities.
    \item Apply one optimizer update.
\end{enumerate}
\end{minipage}
\end{algorithm}
